\makeatletter
\def\input@path{{../styles/}{../../styles/}{../../../styles/}{../}{../../}{../../../}}
\makeatother


\documentclass{ee122_pset}

% Assignment info
\author{\rule{3cm}{0.4pt}} % Name placeholder
\submitdate{\rule{3cm}{0.4pt}} % Submission date placeholder
\problemset{Problem Set \#10: Bode plots for control design}
\renewcommand{\duedate}{April 22, 2026}
\shorttitle{Problem Set \#10}

\begin{document}
\problem{1} \textbf{[40 points]} For a second-order transfer function
\[
G(s) = \frac{\omega_0^2}{s^2 + 2\zeta\omega_0 s + \omega_0^2}
\]
prove that for frequencies that are much higher than the natural frequency $\omega_0$, the magnitude of the transfer function in decibels is approximately $40\log\omega_0 - 40\log\omega$ dB. That is, the magnitude decreases at a rate of 40 dB/decade for frequencies much higher than $\omega_0$. This is the reason why we commonly say that one pole contributes a slope of -20 dB/decade. Here, since there are two poles, the slope is -40 dB/decade.
\problem{2} \textbf{[60 points]} Using MATLAB/Python, plot the Bode magnitude and phase plots for Example 9.13 in the textbook. 
\begin{itemize}
\item Clearly label the frequencies at which the slopes change (10 points)
\item Generate an approximate Bode plot by applying the following logic: a pole at a frequency $\omega_p$ contributes a slope of -20 dB/decade for frequencies higher than $\omega_p$, and a zero at a frequency $\omega_z$ contributes a slope of +20 dB/decade for frequencies higher than $\omega_z$. Remember that a zero of a transfer function is the value of $s$ that makes the transfer function numerator equal to zero. (30 points)
\item Your code should generate same plot as in Figure 9.15 in the textbook. (10 points)
\item Write about your approach in words and include your code in the submission. (10 points)
\end{itemize}

Optional practice problem: Try solving Problem 9.4 --- Using a Bode plot for the analysis of a operational amplifier circuit.
\end{document}